\clearpage

\section{Characterization of Magnetic Phases}\label{sec-nbcp-phase-characterization}

\subsection{Y, UUD, V and the polarized
state}\label{y-uud-v-and-the-polarized-state}

At common azimuth \(\phi=0\), the following signed-angle definitions fix
the states independently of historical names:

\begin{longtable}[]{@{}
  >{\raggedright\arraybackslash}p{(\linewidth - 4\tabcolsep) * \real{0.3333}}
  >{\raggedright\arraybackslash}p{(\linewidth - 4\tabcolsep) * \real{0.3333}}
  >{\raggedright\arraybackslash}p{(\linewidth - 4\tabcolsep) * \real{0.3333}}@{}}
\toprule\noalign{}
\begin{minipage}[b]{\linewidth}\raggedright
State
\end{minipage} & \begin{minipage}[b]{\linewidth}\raggedright
\((\vartheta_A,\vartheta_B,\vartheta_C)\)
\end{minipage} & \begin{minipage}[b]{\linewidth}\raggedright
Characteristic structure
\end{minipage} \\
\midrule\noalign{}
\endfirsthead
\toprule\noalign{}
\begin{minipage}[b]{\linewidth}\raggedright
State
\end{minipage} & \begin{minipage}[b]{\linewidth}\raggedright
\((\vartheta_A,\vartheta_B,\vartheta_C)\)
\end{minipage} & \begin{minipage}[b]{\linewidth}\raggedright
Characteristic structure
\end{minipage} \\
\midrule\noalign{}
\endhead
\bottomrule\noalign{}
\caption{Signed-angle definitions of the three-sublattice reference states.}\label{tab-nbcp-03-phase-diagram-2}\\
\endlastfoot
Y & \((t,-t,\pi)\) & Opposite transverse components; one spin
antiparallel to the field \\
UUD & \((0,0,\pi)\) & Two up spins and one down spin; \(m_z=S/3\) \\
V & \((t,t,v)\) & Two equal canted spins; the third has the opposite
transverse direction \\
P & \((0,0,0)\) & All spins parallel to the field; \(m_z=S\) \\
\end{longtable}

For Y, substituting \(c=\cos t\) into Eq.~\ref{eq-nbcp-three-energy}
gives

\begin{equation}\phantomsection\label{eq-nbcp-y-energy}{
e_Y=S^2[-J+(J+J_z)c^2-2J_zc]-\frac{hS}{3}(2c-1),
\qquad c=\frac{h+3SJ_z}{3S(J+J_z)}.
}\end{equation}

Its merger with UUD occurs at \(c=1\). The collinear energies are
\(e_{\rm UUD}=-S^2J_z-hS/3\) and \(e_P=3S^2J_z-hS\). Their crossing
alone does not locate the intervening canted phase.

For the V family, the energy and two independent stationarity conditions
are

\begin{equation}\phantomsection\label{eq-nbcp-v-stationarity}{
\begin{aligned}
e_V={}&S^2\left[J(\sin^2t+2\sin t\sin v)
+J_z(\cos^2t+2\cos t\cos v)\right]
-\frac{hS}{3}(2\cos t+\cos v),\\
0={}&J\cos t(\sin t+\sin v)-J_z\sin t(\cos t+\cos v)
+\frac{h}{3S}\sin t,\\
0={}&2J\sin t\cos v-2J_z\cos t\sin v+\frac{h}{3S}\sin v.
\end{aligned}
}\end{equation}

The second and third lines are respectively
\((\partial_t e_V)/(2S^2)=0\) and \((\partial_v e_V)/S^2=0\). Their
solutions must also pass a stability check. At \(B=1.4\) T with the
current gap parameters, the selected classical branch has
\((t,t,v)=(0.409364861,0.409364861,-1.223629194)\) rad.

The original {\small\texttt{\nolinkurl{Three_MSL.tex}}} calls a family
with two equal canted spins a ``Gamma phase.'' Its normalized vectors
are

\[
\mathbf n_A=\mathbf n_B=(-\sin t,0,\cos t),
\qquad \mathbf n_C=(\sin\psi,0,-\cos\psi).
\]

A common \(R_z(\pi)\) rotation maps this family to V above with
\(v=\psi-\pi\). This is a geometric identification of the reference
families; at nonzero \(J_\Gamma\), that spin-only rotation is not a
symmetry of the full Hamiltonian. The old ``V'' ansatz instead assigns
opposite transverse components to A and B. Its restricted extrema must
not be substituted for the V solution above or identified with the
paper's \(\Psi\) phase merely by renaming.

\subsection{Competing magnetic cells and quantum
results}\label{competing-magnetic-cells-and-quantum-results}

The source notebook contains a two-sublattice stripe energy and a
four-sublattice variational construction. These are essential
competitors because they retain the anisotropic bond contributions that
cancel in the three-sublattice energy. For a stripe whose
same-sublattice bonds lie along \(\pm\boldsymbol\delta_1\), with
\(\mathsf J_d\) the bond matrix of Eq.~\eqref{eq-nbcp-bond-matrix} for the
family \(\pm\boldsymbol\delta_d\), the bond count gives

\begin{equation}\phantomsection\label{eq-nbcp-stripe-energy}{
E^{(2)}_{\rm cell}=\mathbf S_A^{\mathsf T}\mathsf J_1\mathbf S_A
+\mathbf S_B^{\mathsf T}\mathsf J_1\mathbf S_B
+2\mathbf S_A^{\mathsf T}(\mathsf J_2+\mathsf J_3)\mathbf S_B
-h(S_A^z+S_B^z).
}\end{equation}

Comparison with a three-sublattice state requires
\(E^{(2)}_{\rm cell}/2\) and \(E^{(3)}_{\rm cell}/3\). A stationary
point found in a fixed spin plane still requires fluctuations outside
that plane and at finite momentum to be checked. In particular, an
exactly in-plane stripe at \(h\ne0\) experiences a linear Zeeman torque;
a field-independent Hessian evaluated there does not prove finite-field
stability.

The target paper reports Y, UUD, V, \(\Psi\) and P at zero
bond-dependent SOC, and stripe and four-site SkX orders as SOC grows.
Its phase map uses iDMRG cylinders, while the reported \(\Psi\) state is
not a stable LSWT reference. These results motivate the competitor list;
their phase boundaries and finite-circumference convergence have not
been independently reproduced here. See
Park et al., Section II~\cite{park2026}.

\subsection{Harmonic phase diagram and larger magnetic
cells}\label{sec-nbcp-harmonic-phase-diagram}

The two planes \((h,J_{\rm PD})\) at \(J_\Gamma=0\) and
\((h,J_\Gamma)\) at \(J_{\rm PD}=0\) were scanned at harmonic order with
\(J=0.075\) meV, \(J_z=0.125\) meV, \(S=1/2\) and the field along \(z\)
as the Zeeman energy \(h\), on a grid with field step 0.02 meV
(\href{../../../data-space/verification/261001-nbcp-harmonic-phase-diagram/README.md}{2026-10-01
scan}). At each point the classical minimum on the one-, two-, three-
and four-site cells is tested for LSWT stability on matched momentum
meshes (18 primitive momenta per direction), and the stable states are
ranked by \(E_{\rm cl}+E_{\rm zp}\). For the three-site states, whose
classical energy does not depend on a common rotation about \(z\), twelve
angles in \([0,\pi/3)\) are tested. Global sign changes
\(J_{\rm PD}\to-J_{\rm PD}\) and \(J_\Gamma\to-J_\Gamma\) are removed by
\(R_z(\pi/2)\) and \(R_z(\pi)\), so each axis needs only \(J>0\). On the
PD axis, Y persists to \(J_{\rm PD}\approx0.015\) meV and gives way to a
two-site stripe from about 0.0225 meV; V shrinks from the high-field side
and is replaced above \(h\approx0.36\) meV by the stripe from
\(J_{\rm PD}\approx0.0125\) meV; UUD survives to 0.03 meV and no SkX
appears. On the \(\Gamma\) axis, Y, UUD and V
persist to \(J_\Gamma\approx0.03\) meV, followed by a low-field stripe
and a four-site \(|Q|=2\) SkX. This agrees qualitatively with the
iDMRG phase map of the target paper; its boundary values were not
compared.

The four-site scan leaves two kinds of point undetermined. At 45 points
no candidate is LSWT-stable. At 19 further points the lowest classical
candidate is LSWT-unstable and a stable state of higher classical energy
was ranked first; the classical ground state then lies below an unstable
state and outside the four cells. Both kinds were revisited with
commensurate cells of up to 36 sites (\(3\times3\),
\(2\sqrt3\times2\sqrt3\), \(4\times4\), \(3\sqrt3\times3\sqrt3\) and
\(6\times6\)), twelve random starts per cell together with tiled states
of the smaller cells, so that a larger cell can only lower the energy
(\href{../../../data-space/verification/261005-nbcp-hatched-region/hatched-region-check.json}{cell
search}). The lowest state found is tested for LSWT stability as above;
a state counts as unstable only if it has a negative mode on two meshes
of different parity, because the two-site stripe has an isolated zero
mode at a zone-boundary momentum that one mesh hits exactly. The grid
neighbours of every point whose answer changes were searched in turn
until no new point changed (143 points in all). The results are:

\begin{itemize}
\item
  \emph{Low field between Y and the stripe.} At
  \(J_{\rm PD}=0.0175\)--0.0275 meV and \(h\le0.10\)--0.18 meV
  (34 points) the lowest state is a 12-site
  (\(2\sqrt3\times2\sqrt3\)) state, LSWT-stable. Its longitudinal
  spin component orders at K, as a three-sublattice density, and its
  in-plane component at M, as in a stripe; it lies up to
  \(1.9\times10^{-3}\) meV per spin below the four-site minimum. At
  \(J_{\rm PD}=0.0175\) meV, \(h=0.08\) and 0.10 meV, a 27-site state of
  the same longitudinal order is lower still by
  \((1\)--\(5)\times10^{-6}\) meV. The region extends into the stripe of the
  four-site scan (\(J_{\rm PD}=0.025\) meV, \(h=0.12\)--0.14 meV;
  \(J_{\rm PD}=0.0275\) meV, \(h=0.04\)--0.18 meV), while
  \(J_{\rm PD}=0.03\) meV stays a stripe and the Y and UUD boundaries are
  unchanged. With the zero-point energy on the same meshes, the stripe
  is lower again by \((0.5\)--\(1.2)\times10^{-3}\) meV at
  \(J_{\rm PD}=0.0275\) meV, \(h=0.04\)--0.12 meV; at the other five
  former stripe points the 12-site state also has the lower harmonic
  energy. Where no four-site candidate was stable, the 12-site state is
  the only stable state found. At
  \(J_{\rm PD}=0.015\) meV, \(h=0.02\) meV the three-site state remains
  the lowest and unstable.
\item
  \emph{Near saturation.} On the V side at \(J_{\rm PD}\le0.01\) meV,
  \(h\ge0.40\) meV, and on the \(\Gamma\) axis at
  \(J_\Gamma=0.04\)--0.10 meV, \(h=0.42\)--0.48 meV, no cell up to 36
  sites lowers the three-site state, which stays unstable at 30 points.
  The negative modes, with eigenvalues of \(-3\times10^{-4}\) to
  \(-4\times10^{-3}\) meV, lie mostly at the smallest momenta of the mesh
  and otherwise at single momenta near the zone boundary, which points to
  a long-wavelength or incommensurate modulation. This includes the band
  of four-site states at \(h=0.42\)--0.46 meV and
  \(J_\Gamma=0.04\)--0.075 meV of the four-site scan; the four-site
  states at \(J_\Gamma\ge0.08\) meV, \(h=0.46\) meV and below this band
  remain the lowest and stable. At \(J_{\rm PD}=0.01\)--0.0125 meV,
  \(h=0.42\)--0.50 meV, a 27-site state with in-plane order at a
  wave vector of that cell, e.g.\ \((4/9,4/9)\) in reciprocal-lattice
  units between K and M, lies up to \(7\times10^{-5}\) meV below the
  four-site minimum. It is LSWT-stable at \(J_{\rm PD}=0.01\) meV,
  \(h=0.42\) and 0.44 meV, and at \(J_{\rm PD}=0.0125\) meV,
  \(h=0.48\) meV, with a harmonic energy lower than that of the
  four-site state, and unstable at the other three points. At
  \(J_{\rm PD}=0.0125\) meV, \(h=0.40\)--0.44 meV, the lowest state is
  the two-site stripe, stable on the second mesh; at \(h=0.46\) meV it
  is weakly unstable (\(-1.2\times10^{-5}\) meV).
\item
  \emph{Narrow V windows.} Near the high-field edge of V as few as one
  of the twelve angles is stable. At \(J_\Gamma=0.05\) meV,
  \(h=0.40\) meV the three-site state is now stable at three angles
  (unstable in the four-site scan), and at \(J_\Gamma=0.06\) meV,
  \(h=0.48\) meV it is now unstable at all twelve (stable before). The
  edge of V is therefore resolved only to about one grid step.
\item
  \emph{Reference points.} Y at 0.2~T and V at 1.4~T for
  \(J_{\rm PD}=0.010\) meV remain the lowest states up to 36 sites and
  LSWT-stable, so the starting states of the supersolidity analysis are
  unchanged.
\end{itemize}

These are classical energies and harmonic stability on commensurate
cells of up to 36 sites; an incommensurate modulation is seen only
through its commensurate approximants. The energy differences of order
\(10^{-4}\)--\(10^{-3}\) meV between the 12-site, 27-site, stripe and
four-site states are below the scale of the next order in \(1/S\), and
the zero-point differences rest on coarse meshes for the large cells
whose convergence has not been checked. Points still undetermined (35
of the 143) need a search over continuous wave vectors, such as single-
and double-\(\mathbf Q\) spirals, and a Luttinger--Tisza lower bound.
